\subsection{The actual third obstruction and its integral residue}
\label{sec:ghost-third-obstruction}

\begin{proposition}[The cochain representative of strict tame localization]
\label{prop:ghost-tame-cochain}
Use the morphism \(\Psi_{\ell,r}\) of
Theorem~\ref{thm:ghost-strict-localization}. Define
\(C_t^\bullet\) to be the subcomplex of normalized
continuous global cochains with values in
\(\mathbb Z_\ell(r)\) whose restrictions, as actual
cochains, at every \(q\ne\ell\) factor through
\[
 G_q^{t,\ell}=\mathbb Z_\ell(1)\rtimes\widehat{\mathbb Z},
       \qquad FT F^{-1}=T^q,
\]
where the second factor is generated by arithmetic Frobenius.
The coefficient action factors through these quotients,
so the differential preserves this subcomplex. The
restriction of \(\Psi_{\ell,r}\) along the natural map
\(H^2(C_t^\bullet)\to H^2_{\rm cts}
(G_{\mathbb Q},\mathbb Z_\ell(r))\) is induced by the
following cochain map in degrees one and two:
\[
 b\longmapsto (b(T_q))_{q\ne\ell},\qquad
 c\longmapsto (z_q(c))_{q\ne\ell},
\]
with target \([R_\ell\xrightarrow{d_r}R_\ell][-1]\),
where
\[
 z_q(c)=q^{-1}\left(c(F,T)-c(T^q,F)
                         -\sum_{j=1}^{q-1}c(T^j,T)\right).
\]
In particular the representative uses the literal global
cochain and retains its derived-limit contribution.
\end{proposition}
\begin{proof}
At \(q\ne\ell\), inertia acts trivially on the coefficient
module, whereas Frobenius acts by \(q^r\). For a normalized
one-cochain \(b\), direct substitution of
\(\delta b(g,h)=g b(h)-b(gh)+b(g)\), using \(FT=T^qF\),
gives
\[
 qz_q(\delta b)=(q^r-q)b(T),\qquad
       z_q(\delta b)=(q^{r-1}-1)b(T).
\]
Indeed the sum of \(\delta b(T^j,T)\) telescopes to
\(qb(T)-b(T^q)\). Thus in degrees one and two the
displayed formula is the cochain map
\[
 b\longmapsto b(T),\qquad c\longmapsto z_q(c)
\]
to \([T_n(-1)\xrightarrow{q^{r-1}-1}T_n(-1)][-1]\)
after any finite coefficient reduction. It vanishes on
normalized unramified cochains.

We check its identification with the supports morphism,
including the unit \(q^{-1}\). In the extension defined
by \(c\), let \(\tau=(0,T)\) and \(\phi=(0,F)\).
The central defect of \(\phi\tau=\tau^q\phi\) is
exactly \(qz_q(c)\), by multiplication in the extension:
\(\tau^q\) has central coordinate
\(\sum_{j=1}^{q-1}c(T^j,T)\). The extension over tame
inertia splits by the powers of a lift of its generator.
Frobenius acts on an inertia homomorphism \(u\) by
\[
       (F\cdot u)(T)=F\,u(F^{-1}TF)=q^{r-1}u(T).
\]
The discrepancy between this conjugated inertia splitting
and the original one, evaluated at \(T\), is therefore
\(q^{-1}\) times the displayed relator, namely \(z_q(c)\).
This is the degree-two edge map obtained from the two
successive cyclic resolutions for inertia and Frobenius.
Its degree-one map is restriction to inertia. These are
precisely the relative cochain maps for unramified
cochains inside tame cochains. Purity and the
Hochschild--Serre calculation identify that relative
complex with \(R\Gamma(\mathbb F_q,T_n(-1))[-1]\),
as in the proof of Theorem~\ref{thm:ghost-strict-localization}.
Wild inertia and the remaining prime-to-\(\ell\) tame
inertia have exact invariant functors on these coefficients:
normalized averaging over each finite quotient has
denominator prime to \(\ell\). This gives the same
identification for the local Galois group. Every map just
described commutes with coefficient reduction.

Finally \((z_q(c))\) belongs to \(R_\ell\). For each
\(n\), the global cochain modulo \(\ell^n\) factors
through a finite Galois quotient. Outside its finite
ramification set, \(T\) and every \(T^j\) have trivial
image; normalization makes every term of the formula
zero modulo \(\ell^n\). Division by \(q\) is by an
\(\ell\)-adic unit. Thus the coefficient-compatible
maps assemble into the literal inverse-limit morphism
defining \(\Psi_{\ell,r}\), and give the claimed
representative. The same argument for \(b(T)\) shows
that a global coboundary changes it by an element of
\(d_rR_\ell\), with no closure taken.
\end{proof}

\begin{proposition}[Complete integral test for a specified class-two section]
\label{prop:ghost-class3-residue}
Use \(x,y,C=[x,y],D=[C,x],F_3=[C,y]\), with
\([u,v]=uvu^{-1}v^{-1}\), and the tangential splitting of
Corollary~\ref{cor:ghost-class2-lift}. Write a specified
class-two section as
\(s(g)=y^{b(g)}x^{a(g)}C^{h(g)}s_0(g)\).
Set \(\epsilon=(\chi-1)/2\) and
\(\eta=(\chi^2-1)/24\). Its obstruction to the actual
class-three arithmetic quotient has the two continuous
cochain representatives
\[
 \begin{split}
 W_D&=-(a+\epsilon)\cup h-\binom a2\cup b,\\
 W_F&=(b+\epsilon)\cup(ab-h)+\binom b2\cup a-\eta\cup b.
 \end{split}
\]
The pointwise products and binomial expressions in this
formula are weight-two cochains. For \(q\ne\ell\), put
\[
 \begin{gathered}
 (m,n,t)=(a(T),b(T),h(T)),\qquad
 (\alpha,\beta,\gamma)=(a(F),b(F),h(F)),\\
 e=(q-1)/2,\qquad e_2=(q^2-1)/24,
 \end{gathered}
\]
and define the \(\ell\)-adically integral expression
\[
 B_q(z)=q^{-1}\sum_{j=1}^{q-1}\binom{jz}{2}
       =\frac{q-1}{12}\bigl((2q-1)z^2-3z\bigr).
\]
The images of the two obstruction classes under
\(\Psi_{\ell,3}\) are represented by the actual families
\[
 \begin{split}
 Z_{D,q}={}&m\gamma-q(\alpha+e)t-n\binom\alpha2
       +q^{-1}\beta\binom{qm}{2}+emt+nB_q(m),\\
 Z_{F,q}={}&n(\gamma-\alpha\beta)+q(\beta+e)(mn-t)
       +m\binom\beta2-e_2n-q^{-1}\alpha\binom{qn}{2}\\
       &\hspace{10mm}-e(mn^2-nt)-mB_q(n).
 \end{split}
\]
Both families belong to \(R_\ell\). The section \(s\)
lifts to the full integral class-three quotient exactly
when, for every \(\ell\), both coordinatewise quotients
\[
       \left(\frac{Z_{D,q}}{q^2-1}\right)_{q\ne\ell},
       \qquad
       \left(\frac{Z_{F,q}}{q^2-1}\right)_{q\ne\ell}
\]
are integral and tend to zero \(\ell\)-adically.
This confirms SECTIONS A.43.1 and A.46.5--A.46.6 with
their actual continuous obstruction classes.
\end{proposition}
\begin{proof}
Discrete collection through weight three gives
\[
 x^a y^b=y^b x^a C^{ab}
                  D^{b\binom{a+1}{2}}F_3^{a\binom{b+1}{2}}.
\]
The weight-two coefficient of the actual path cocycle
\(\varphi_g\) is \(\eta(g)\); the precise path and
normalization are given by
\cite[(12.12), (12.19), (12.35)]{GhostWickelgren2012}.
Its integrality includes two and three, since
\(\chi^2-1\in24\widehat{\mathbb Z}\).
Consequently the genuine arithmetic action through
weight three is
\[
 g(y^b x^a C^h)=
 y^{\chi b}x^{\chi a}C^{\chi^2h}
 D^{-\epsilon\chi^2h}
 F_3^{-\epsilon\chi^2h-\eta(g)\chi b}.
\]
Keep the first three Hall coordinates and set the final
two to zero to obtain an actual continuous set section.
Its defect, using the displayed collection and action,
has coordinates
\[
 \begin{split}
 c_D(g,k)={}&h(g)\chi(g)a(k)
   +\binom{a(g)+1}{2}\chi(g)b(k)
   +a(g)\chi(g)^2a(k)b(k)-\epsilon(g)\chi(g)^2h(k),\\
 c_F(g,k)={}&h(g)\chi(g)b(k)
   +a(g)\binom{\chi(g)b(k)+1}{2}
   -\epsilon(g)\chi(g)^2h(k)-\eta(g)\chi(g)b(k).
 \end{split}
\]
They are normalized two-cocycles in the central module
\(\widehat{\mathbb Z}(3)^2\). The class-two cocycle
equation is \(\delta h=-a\cup b\). Direct expansion
now gives the exact coboundary comparisons
\[
 W_D=c_D+\delta(ha),\qquad
 W_F=c_F+\delta\left(hb-a\binom b2\right).
\]
These are integral continuous cochains. This also
independently supplies the comparison underlying
\cite[Propositions 13 and 17]{GhostWickelgren2012},
with the two abelian coordinate names interchanged to
match the present convention.

For \(q\ne\ell\) the action through class three is
unramified, by the displayed action formula and the
unramifiedness of \(\chi_\ell\). A local class-two
section kills wild and non-\(\ell\) tame inertia:
inertia acts trivially, so its geometric restriction is
a homomorphism from those prime-to-\(\ell\) groups to
a pro-\(\ell\) group. Its Hall coordinate cochains,
and hence \(W_D,W_F\), factor through \(G_q^{t,\ell}\).
Thus Proposition~\ref{prop:ghost-tame-cochain} applies.
Substituting \(a(T^j)=jm\), \(b(T^j)=jn\),
\(\chi(T^j)=1\) into that formula gives exactly
\(Z_D,Z_F\). The displayed power sum computes \(B_q\)
and proves its integrality, including \(\ell=2,3\),
without dividing by a nonunit. All identities are
also verified symbolically in
\path{code/check_class3_residue.py}; its output is retained.

Finally the obstruction vanishes exactly when both
central cochains admit continuous primitives. Such
primitives correct the last two coordinates of the
set section and give precisely a section of the actual
class-three quotient. By
Theorem~\ref{thm:ghost-strict-localization},
\(\Psi_{\ell,3}\) is an isomorphism; its zero class
means literal membership in \(d_3R_\ell\). Multiplication
by each \(q^2-1\) is injective on \(\mathbb Z_\ell\),
so this is exactly the integrality and convergence
condition stated. Products over \(\ell\) give the
full profinite central layer, with continuous primitives
assembled componentwise.
\end{proof}

\begin{proposition}[What the prime \(1129\) does to the third obstruction]
\label{prop:ghost-1129-unramified}
For the ordinary factorial witness of
Theorem~\ref{thm:ghost-factorial-selmer32}, every global
class-two section has zero local class-three obstruction
at \(1129\), for every coefficient prime.
At every coefficient prime \(\ell\ne1129\), its
restriction there lifts to the entire geometrically
pro-\(\ell\) local arithmetic fundamental group.
\end{proposition}
\begin{proof}
The prescribed ordinary local point is a unit and its
complement is a unit at \(q=1129\). Thus its two
integral abelian inertia coordinates are zero. For
\(\ell\ne q\), the remaining class-two inertia is a
homomorphism into the central torsion-free module
\(\mathbb Z_\ell(2)\). Frobenius conjugation and the
tame relation give \((q^2-q)h(T)=0\), hence \(h(T)=0\).
Wild and non-\(\ell\) tame inertia have zero image as
well. The given local section is therefore unramified.

The arithmetic action on the entire free pro-\(\ell\)
geometric group is unramified at \(q\). Here is a
proof directly from its nilpotent quotients. The
class-two action is the pure cyclotomic action already
computed. Inductively, the kernel of the automorphism
map from class \(i+1\) to class \(i\) is a subgroup
of \(\operatorname{Hom}(N_1,L_{i+1})\), a torsion-free
\(\mathbb Z_\ell\)-module on which Frobenius acts by
\(q^i\). An inertia image in this kernel is killed by
\(q^i-q\), by the same tame relation; it is zero for
\(i\ge2\). Wild and prime-to-\(\ell\) inertia also
vanish in that module. Thus inertia acts trivially at
every class. The free pro-\(\ell\) group is the inverse
limit of these quotients: every finite \(\ell\)-group
has finite nilpotency class. Hence inertia acts trivially
on the full group.

The local unramified arithmetic group is consequently
\(\Pi^{(\ell)}\rtimes\widehat{\mathbb Z}\), where
Frobenius acts by the actual arithmetic automorphism.
Choose any lift of the given section's Frobenius element
from class two to this group. Taking its powers gives
a continuous homomorphism from
\(\widehat{\mathbb Z}\): a single element in a profinite
group defines such a homomorphism by the universal
property of the procyclic completion. Its composition
with the arithmetic projection sends Frobenius to
Frobenius, so is the identity. Inflation to the local
Galois group gives the required section, reducing
exactly to the specified class-two section.

For \(\ell=q=1129\), local Tate duality identifies the
dual of \(H^2(G_{\mathbb Q_q},\mathbb Z_q(3))\)
with \((\mathbb Q_q/\mathbb Z_q)(-2)^{G_{\mathbb Q_q}}\).
The cyclotomic image modulo \(q\) has order \(q-1\),
which does not divide two. Some cyclotomic scalar
therefore acts with difference from one a \(q\)-adic
unit, making these invariants zero
\cite[I.2.3]{GhostMilneADT}. The local third obstruction
vanishes also at this coefficient prime.

Thus the finite Selmer exclusion at \(1129\) does not
become an unrestricted third obstruction through local
restriction. The exact local cup-residue map multiplies
the central comparison by the two valuation coordinates,
both of which are zero here; the unramified lift above
exhibits this zero directly at every higher
prime-to-\(1129\) stage.
\end{proof}
